\documentclass{article}

\usepackage{hyperref}

\title{CS536 Final Project}
\author{Author}
\date{}

\begin{document}
\maketitle

\section{Motivation}
Throughout this course, we have mostly considered learning settings where there is a single entity handling training and testing.
However, these techniques are not always applicable to scenarios like edge computing and IoT devices.
Often, there may be multiple devices with disjoint data sets that need to be aggregated to train a shared model.
These devices may also have bandwidth or privacy requirements, preventing users from directly aggregating data then training.
This is addressed by techniques in federated learning \cite{yang2019}.

There are several ways data can be split among multiple devices.
In the first case, each device can hold a subset of the data, recording different events with identical features.
This is the problem that \underline{horizontal federated learning} solves \cite{li2020}.
Alternatively, perhaps each device records some information about the same event, but they hold a disjoint set of features.
We address this problem with \underline{vertical federated learning} \cite{feng2020}.


\bibliographystyle{plain}
\bibliography{report.bib}

\end{document}